\documentclass[12pt, ukrainian]{article}
\setlength\textwidth{190mm} \setlength\hoffset{-25mm}
\setlength\textheight{277mm} \setlength\voffset{-28mm}
%\setlength\textwidth{190mm}%\setlength\voffset{-38mm} \setlength\hoffset{-38mm}
%\setlength\parindent{0mm}
\pagestyle{empty}
\usepackage[T2A]{fontenc}
\usepackage[utf8]{inputenc}
\usepackage[ukrainian]{babel}
\usepackage{graphicx}
\usepackage{caption}
\usepackage{verbatim, listings, clrscode3e}
\usepackage{fancyvrb}
\DeclareCaptionFormat{nocaption}{}
\captionsetup{format=nocaption,aboveskip=0pt,belowskip=0pt}
\usepackage[Export]{adjustbox}
\adjustboxset{max size={0.9\linewidth}{0.9\paperheight}}
\def\code#1{\Verb!#1!}
\begin{document}

{{\begin {small}\par
{ \center  Київський національний університет імені Тараса Шевченка \par}
{\parbox {5 cm}{Освітня програма \hfill}{\parbox {8 cm}{Прикладна математика, Системний аналіз \hfill}}\parbox {5 cm}{\hfill Семестр 2}}\par
{\parbox {5 cm} {Навчальний предмет \hfill}\parbox {7 cm}{Програмування \hfill}\parbox {6cm}{\hfill}\par}\vspace{-7pt} \end{small}}{\center Контрольна робота, варіант  \No \textbf
{ 20 }\par}}
{
\begin {large}
{\begin {center}\textbf{Завдання 1}\end {center}}\par
Що буде надруковано за виконання?\par\par
\begin{center}
	\adjustimage{max size={0.9\linewidth}{0.9\paperheight}}
	{../pict/20_1.png}
\end{center}
\newpage

{\begin {center}\textbf{Завдання 2}\end {center}}\par
Нехай змінна \kw{s} позначає послідовність цілих, по якій можна ітерувати.

Написати вираз, що перевіряє, чи правда, що послідовність  містить число 13.\par
\newpage

{\begin {center}\textbf{Завдання 3}\end {center}}\par
Що буде надруковано за виконання?
%Оригинал был утерян. Требует отладки!\par
\begin{center}
	\adjustimage{max size={0.9\linewidth}{0.9\paperheight}}
	{../pict/20_3.png}
\end{center}
\newpage

{\begin {center}\textbf{Завдання 4}\end {center}}\par
Що буде надруковано за виконання?\par
\begin{center}
	\adjustimage{max size={0.9\linewidth}{0.9\paperheight}}
	{../pict/20_4.png}
\end{center}
\newpage

{\begin {center}\textbf{Завдання 5}\end {center}}\par
Змінні \kw{next} та \kw{prev} вузла списку позначають наступний та попередній за ним вузол відповідно. Починаючи з голови, у список записано послідовні цілі числа від~$-21$ до~$45$. Нехай змінна \kw{p} позначає вузол, що зберігає значення~$6$.
Яке значення зберігається у вузлі \kw{p.next.next.prev.prev.next} ?\par
\newpage

{\begin {center}\textbf{Завдання 6}\end {center}}\par

\begin{center}
	\adjustimage{max size={0.8\linewidth}{0.8\paperheight}}
	{../BinTree/trees_exp/72.png}
\end{center}
Указати послідовність міток вершин дерева, зображеного на рис., що утвориться в результаті обходу дерева в оберненому порядку. Мітки записувати через пробіл.\par Алгоритм починає роботу з кореня (на рис. це верхня вершина).\par
\newpage

{\begin {center}\textbf{Завдання 7}\end {center}}\par
Для графа на рисунку вказати послідовність, в якій алгоритм пошуку в глибину буде відмічати відвідані вершини, якщо списки суміжності вершин переглядаються алгоритмом за зростанням міток вершин та вершини відмічаються на першому вход\-женні у вершину.\par
Пошук розпочинається з вершини 2.\par

\begin{center}
	\adjustimage{max size={0.9\linewidth}{0.9\paperheight}}
	{../Graphs/Traversals/178.png}
\end{center}\par
\newpage

{\begin {center}\textbf{Завдання 8}\end {center}}\par
Для графа на рисунку з вершини 1 запущено алгоритм пошуку в ширину, що будує кістякове дерево. Списки суміжності вершин переглядаються алгоритмом за зростанням міток вершин.\par
\begin{center}
	\adjustimage{max size={0.9\linewidth}{0.9\paperheight}}
	{../Graphs/Traversals/276.png}
\end{center}
\par Які з наведених ребер входять до складу отриманого кістякового дерева?\par
1) (1, 3)\par
2) (0, 6)\par
3) (0, 4)\par
4) (0, 2)\par
5) (3, 4)\par
6) (1, 2)\par
{\vskip 15pt} У формі вибрати номери відповідних ребер.


Якщо заповнюєте текстовий файл, а не гугл-форму, то потрібно через пробіл записа\-ти вибрані номери.\par
\newpage


Файл \code{'1.txt'} містить журнал з рядками, в яких через кому записано поряд\-ко\-вий номер запису в журналі (ціле), тип події (ERROR, WARNING, INFO, STATUS), ім'я користувача (складається з латиниці), час події (фор\-мат як у прикладі), код події (ціле).

Білі символи навколо роздільників не допускаються.

Приклад рядка файлу



\begin{verbatim}
13,ERROR,username,2022-05-05 13:26,404
\end{verbatim}


Нижче наведено код, який шукає усі коди помилок (подія ERROR), що зустріча\-ють\-ся в журналі, та зберігає їх у змінній \kw{codes}.



\begin{verbatim}
codes = set()
regex = re.compile(r'XXX')
with open('1.txt', encoding='utf8') as f:
    for s in f:
        res = re.fullmatch(regex, s)
        if not res: continue
        s = ''
        codes.add(r'YYY')
\end{verbatim}


{\begin {center}\textbf{Завдання 9}\end {center}}\par
Який рядок треба записати на місце \kw{XXX} ?\par
\vspace{5mm}


{\begin {center}\textbf{Завдання 10}\end {center}}\par
Який рядок треба записати на місце \kw{YYY} ?\par


\newpage


\end {large}}

\end{document}
